\FloatBarrier
\section{Decisiones de diseño}\label{sec:diseno}

Cada vez que el equipo tomó una decisión, la anotamos con un número (D-01, D-02 y así) en el archivo \tt{docs/llm/decisions.md} del repositorio. Cuando una duda no tenía respuesta, la dejamos como pregunta abierta y la resolvimos con el equipo, nunca en silencio. Esta sección explica, en palabras simples, las decisiones que más cambiaron el sistema.

\subsection{Las herramientas}\label{sec:pila}

El equipo no tiene un servidor propio, así que elegimos servicios en la nube con plan gratuito. Usamos Next.js, que junta las pantallas y el servidor en un solo proyecto \parencite{nextjs}. La aplicación se publica en Vercel y los datos viven en una base PostgreSQL en Neon (D-01, D-02). Elegimos esto porque un solo proyecto es más fácil de entender para todo el equipo, y porque Vercel crea una versión de prueba de la aplicación para cada cambio, así cualquiera podía probarlo antes de unirlo. El primer prototipo que hicimos quedó solo como referencia (D-05).

\subsection{El código separado en capas}\label{sec:capas}

La entrega 1 pedía que el sistema no fuera un solo bloque donde todo depende de todo. Por eso separamos el código en cuatro capas, cada una en su carpeta, como muestra la \fref{fig:capas}. Una capa solo habla con la de abajo, así que una pantalla nunca toca la base de datos directo. Las reglas del negocio, como cuándo se puede emitir un certificado, viven aparte en \tt{src/lib/domain/}. Como esas reglas no usan la base de datos, se pueden probar en milisegundos.

\begin{figure}[H]
  \centering
  % Layers of the app. A request goes down one layer at a time.
% Redrawn from the development view in docs/diseno-arquitectonico-es.md.
\begin{tikzpicture}[node distance=0.55cm]
  \tikzset{capa/.style={bloque, minimum width=6.4cm, text width=6.0cm}}
  \node[capa] (pant) {\textbf{Pantalla}\\\scriptsize lo que ve el usuario, en español (\tt{page.tsx})};
  \node[capa, below=of pant] (acc) {\textbf{Server action o ruta de la API}\\\scriptsize revisa el rol y lee el formulario (\tt{actions.ts}, \tt{/api/v1})};
  \node[capa, below=of acc] (uc) {\textbf{Caso de uso y reglas del negocio}\\\scriptsize funciones puras con pruebas (\tt{src/lib/domain}, \tt{src/lib/api})};
  \node[capa, below=of uc] (dat) {\textbf{Acceso a datos}\\\scriptsize consultas con Drizzle (\tt{data.ts}, \tt{schema.ts})};
  \node[bloque, below=of dat, minimum width=6.4cm, text width=6.0cm, fill=acento!12] (db) {\textbf{Base de datos}\\\scriptsize PostgreSQL en Neon};

  \node[bloque, right=1.3cm of uc, text width=3.0cm] (srv) {\textbf{Servicios}\\\scriptsize PDF, correo con Resend, exportación a Excel};

  \draw[flecha] (pant) -- (acc);
  \draw[flecha] (acc) -- (uc);
  \draw[flecha] (uc) -- (dat);
  \draw[flecha] (dat) -- (db);
  \draw[flecha] (uc) -- (srv);

  \node[etiq, left=0.4cm of acc, text width=2.6cm, align=right] {si el rol no tiene permiso, aquí se detiene};
  \node[etiq, left=0.4cm of dat, text width=2.6cm, align=right] {el cambio y su fila de bitácora se guardan juntos};
\end{tikzpicture}

  \caption{Las capas de la aplicación y el camino de una petición}\label{fig:capas}
  \backline{fig:capas}
\end{figure}

\subsection{Los datos}\label{sec:datos}

Los límites de calidad pueden ser internacionales o propios de un cliente. La entrega 1 guardaba los internacionales como si fueran de un cliente falso con la clave 999. Nosotros usamos dos tablas separadas, una para los límites internacionales y otra para los de cada cliente (D-07). Un cliente sin límites propios usa los internacionales. Así ninguna lista ni reporte tiene que acordarse de esconder al cliente falso.

Otras dos decisiones protegen la historia. Una inspección guarda los resultados de las dos máquinas juntos (D-09), y el certificado guarda su propia copia de cada valor con su límite (D-10). Si mañana cambian los límites, un certificado viejo sigue diciendo exactamente lo mismo. La \fref{fig:datos} muestra las tablas principales, y el diagrama completo está en \tt{docs/diseno-base-datos-es.md}.

\begin{figure}[H]
  \centering
  \resizebox{\linewidth}{!}{% Data model reduced to the main tables and their relations.
% Redrawn from the conceptual model in docs/diseno-base-datos-es.md. Column lists are cut.
% Each arrow goes from the "one" side to the "many" side of the relation.
\begin{tikzpicture}
  \tikzset{t/.style={tabla, text width=3.2cm}}
  \node[t] (clients)  at (0,0)      {\textbf{clients}\nodepart{two}clave SAP, RFC único,\\contacto, estatus};
  \node[t] (climits)  at (5.4,0)    {\textbf{client\_limits}\nodepart{two}límite propio de un\\cliente por parámetro};
  \node[t] (ilimits)  at (10.8,0)    {\textbf{international\_limits}\nodepart{two}un límite por\\parámetro};

  \node[t] (certs)    at (0,-3)     {\textbf{certificates}\nodepart{two}folio, estatus, envío,\\hash, cancelación};
  \node[t] (lots)     at (5.4,-3)   {\textbf{lots}\nodepart{two}producto, fechas de\\producción y caducidad};
  \node[t] (params)   at (10.8,-3)   {\textbf{parameters}\nodepart{two}los diez factores\\que se miden};

  \node[t] (lines)    at (0,-6)     {\textbf{certificate\_lines}\nodepart{two}copia del valor, del\\límite y la desviación};
  \node[t] (insp)     at (5.4,-6)   {\textbf{inspections}\nodepart{two}secuencia A, B, C,\\motivo de reinspección};
  \node[t] (res)      at (10.8,-6)   {\textbf{inspection\_results}\nodepart{two}un valor por\\parámetro};

  \node[t] (deliv)    at (0,-9)     {\textbf{deliveries}\nodepart{two}cada intento de\\envío por correo};
  \node[t] (equip)    at (10.8,-9)   {\textbf{equipment}\nodepart{two}clave generada,\\tipo, estatus};

  \draw[flecha] (clients) -- node[etiq, above] {tiene} (climits);
  \draw[flecha] (params) -- node[etiq, above, sloped] {acota} (climits);
  \draw[flecha] (params) -- node[etiq, right] {uno a uno} (ilimits);
  \draw[flecha] (clients) -- node[etiq, left] {recibe} (certs);
  \draw[flecha] (lots) -- node[etiq, above] {se certifica en} (certs);
  \draw[flecha] (lots) -- node[etiq, right] {tiene} (insp);
  \draw[flecha] (insp) -- node[etiq, above] {contiene} (res);
  \draw[flecha] (params) -- node[etiq, right] {es de} (res);
  \draw[flecha] (equip) -- node[etiq, right] {mide} (res);
  \draw[flecha] (insp) -- node[etiq, below, sloped] {respalda} (certs);
  \draw[flecha] (certs) -- node[etiq, left] {imprime 10} (lines);
  \draw[flecha] (certs.west) -- ++(-0.5,0) |- node[etiq, pos=0.25, left] {se envía en} (deliv.west);
\end{tikzpicture}
}
  \caption{Las tablas principales de la base de datos y sus relaciones}\label{fig:datos}
  \backline{fig:datos}
\end{figure}

\subsection{Quién puede hacer qué}\label{sec:permisos}

Esconder un botón en la pantalla no protege nada, porque alguien puede escribir la dirección de la página a mano. Por eso el servidor revisa el tipo de cuenta antes de cada acción y de cada consulta. Las cuentas las crea el administrador y no existe el registro público (D-14). Además, el inicio de sesión se bloquea por un minuto después de cinco intentos fallidos (D-28).

\subsection{Las reglas del certificado}\label{sec:reglas}

Un certificado no debe mentir sobre la calidad del lote. Para eso el sistema sigue estas reglas.

\begin{itemize}
  \item El certificado usa siempre la última inspección del lote. Si está incompleta, no se puede emitir, aunque exista otra anterior completa (D-27). Esta regla salió de una revisión del profesor.
  \item Un valor fuera de rango no impide emitir, pero se ve marcado en la columna de diferencia (BR-CER-04).
  \item Antes de emitir se piden la factura, la cantidad, la fecha de envío, el destino y el transportista (BR-CER-05, D-24).
  \item Al emitir, el sistema calcula una huella digital del contenido (D-26). Es un código que cambia si cambia un solo dato, así que sirve para saber si alguien alteró el certificado.
  \item Un certificado emitido nunca se borra. Si tiene un error, se cancela con un motivo y se emite uno nuevo (D-25).
  \item Repetir la inspección de un lote ya completo solo se permite por una falla del equipo, con el motivo escrito (D-15).
\end{itemize}

Estas reglas quitan el riesgo de \enquote{ajustar} una inspección para un cliente. La \fref{fig:certificado} muestra los estados por los que pasa un certificado.

\begin{figure}[H]
  \centering
  % Certificate life cycle. Redrawn from the certificate state diagram in
% docs/especificacion-tecnica-es.md (draft, issued, sent, cancelled).
\begin{tikzpicture}[node distance=1.2cm and 2.4cm]
  \node[estado] (draft) {Borrador};
  \node[estado, right=of draft] (issued) {Emitido};
  \node[estado, right=of issued] (sent) {Enviado};
  \node[estado, below=1.8cm of $(issued)!0.5!(sent)$, fill=white] (canc) {Cancelado};

  \draw[flecha] ($(draft.west)+(-1.0,0)$) -- node[etiq, above] {crear} (draft);
  \draw[flecha] (draft) -- node[etiq, above] {emitir sin\\bloqueos} (issued);
  \draw[flecha] (issued) -- node[etiq, above] {envío\\exitoso} (sent);
  \draw[flecha] (issued) -- node[etiq, left] {cancelar\\con motivo} (canc);
  \draw[flecha] (sent) -- node[etiq, right] {cancelar\\con motivo} (canc);

  \draw[flecha] (draft) to[out=60, in=120, looseness=4] node[etiq, above] {editar} (draft);
  \draw[flecha] (issued) to[out=60, in=120, looseness=4] node[etiq, above] {envío fallido, queda la evidencia} (issued);
  \draw[flecha] (sent) to[out=60, in=120, looseness=4] node[etiq, above] {reenviar} (sent);

  \node[etiq, below=0.35cm of draft, text width=3.2cm] {un borrador no se cancela, solo no se emite};
\end{tikzpicture}

  \caption{Los estados de un certificado}\label{fig:certificado}
  \backline{fig:certificado}
\end{figure}

\subsection{La bitácora}\label{sec:bitacora}

El curso pide saber quién hizo cada acción, con su número de empleado y su nombre (RF-01). El sistema guarda cada cambio y su renglón de bitácora en una sola operación, así que o se guardan los dos o no se guarda ninguno. Si se guardaran por separado, una falla a la mitad podría dejar un cambio sin rastro. La bitácora solo recibe renglones nuevos y nunca se edita ni se borra.

\subsection{Lo que decidimos no construir}\label{sec:fuera}

Algunas ideas quedaron fuera a propósito, porque no cabían en el plazo o porque el riesgo era mayor que el beneficio.

\begin{itemize}
  \item Límites distintos por cliente y por producto. Por ahora los límites son solo por cliente (D-12).
  \item Una firma electrónica con criptografía. La firma imprime el nombre, el número de empleado y la fecha de quien emite (D-18).
  \item La conexión con SAP. Los clientes se capturan a mano en esta versión (RF-03).
  \item La humedad como undécimo valor. Se quedan los diez valores actuales, con rangos de demostración (D-30 a D-32).
  \item Un catálogo de productos. El producto de un lote es texto libre con sugerencias (D-21).
\end{itemize}
